% ch04_fonctions.tex — Fonctions et Port\'ee
\chapter{Fonctions et Port\'ee}
\label{ch:fonctions}

\begin{intuition}
Une fonction est comme une \emph{machine} dans un laboratoire~: on y introduit des
\emph{entr\'ees} (r\'eactifs), elle effectue un \emph{traitement} (r\'eaction chimique) et
produit une \emph{sortie} (produit). Plut\^ot que de reconstruire la machine \`a chaque
exp\'erience, on la r\'eutilise. C'est le principe \textbf{DRY}\index{DRY}~:
<<~\emph{Don't Repeat Yourself}~>>.
\end{intuition}

% ============================================================
\section{Pourquoi utiliser des fonctions~?}
\index{fonction}

\begin{itemize}
    \item \textbf{R\'eutilisabilit\'e}~: \'ecrire une formule une seule fois et l'appeler partout.
    \item \textbf{Lisibilit\'e}~: un programme d\'ecoup\'e en fonctions est plus facile \`a comprendre.
    \item \textbf{D\'ebogage}~: on peut tester chaque fonction ind\'ependamment.
    \item \textbf{Modularit\'e}\index{modularit\'e}~: chaque fonction a une responsabilit\'e unique.
\end{itemize}

% --- TikZ machine ---
\begin{figure}[ht]
\centering
\begin{tikzpicture}[
    box/.style={rectangle, draw, fill=blue!10, rounded corners,
        minimum width=3cm, minimum height=1.5cm, align=center, font=\bfseries},
    io/.style={rectangle, draw, fill=orange!15, rounded corners,
        minimum width=2cm, minimum height=0.8cm, align=center},
    arr/.style={->, thick, >=stealth}
]
\node[io] (in) {Entr\'ees\\(param\`etres)};
\node[box, right=2cm of in] (fn) {Fonction\\(traitement)};
\node[io, right=2cm of fn] (out) {Sortie\\(\py{return})};
\draw[arr] (in) -- (fn);
\draw[arr] (fn) -- (out);
\end{tikzpicture}
\caption{Une fonction vue comme une machine~: entr\'ees $\to$ traitement $\to$ sortie.}
\label{fig:fonction-machine}
\end{figure}

% ============================================================
\section{D\'efinir et appeler une fonction}
\index{def}\index{return}

\begin{definition}[Syntaxe d'une fonction]
\begin{minted}{python}
def nom_fonction(parametre1, parametre2):
    """Docstring : description de la fonction."""
    # corps de la fonction
    return resultat
\end{minted}
\end{definition}

\begin{codeblock}[title=\textbf{Python}]
\begin{minted}{python}
def energie_cinetique(masse, vitesse):
    """Calcule l'énergie cinétique E = 0.5 * m * v²."""
    return 0.5 * masse * vitesse**2

# Appel de la fonction
E = energie_cinetique(2.0, 3.0)
print(f"E = {E} J")
\end{minted}
\end{codeblock}

\begin{codeoutput}
\begin{verbatim}
E = 9.0 J
\end{verbatim}
\end{codeoutput}

% ============================================================
\section{Docstrings}
\index{docstring}

Une \emph{docstring} est une cha\^ine de documentation plac\'ee juste apr\`es \py{def}.
Elle est accessible via \py{help(nom_fonction)}.

\begin{codeblock}[title=\textbf{Python}]
\begin{minted}{python}
def force_gravitationnelle(m1, m2, r):
    """
    Calcule la force gravitationnelle entre deux masses.

    Paramètres
    ----------
    m1 : float — masse 1 (kg)
    m2 : float — masse 2 (kg)
    r  : float — distance entre les centres (m)

    Retour
    ------
    float — force en Newtons
    """
    G = 6.674e-11  # constante gravitationnelle
    return G * m1 * m2 / r**2

# Force Terre-Lune
F = force_gravitationnelle(5.972e24, 7.342e22, 3.844e8)
print(f"Force Terre-Lune : {F:.2e} N")
\end{minted}
\end{codeblock}

\begin{codeoutput}
\begin{verbatim}
Force Terre-Lune : 1.98e+20 N
\end{verbatim}
\end{codeoutput}

% ============================================================
\section{Arguments par d\'efaut et arguments nomm\'es}
\index{argument par d\'efaut}\index{argument nomm\'e}

\begin{codeblock}[title=\textbf{Python}]
\begin{minted}{python}
def chute_libre(t, g=9.81, v0=0):
    """Position d'un objet en chute libre : y = v0*t + 0.5*g*t²."""
    return v0 * t + 0.5 * g * t**2

# Utilisation avec valeurs par défaut
print(f"Terre : y = {chute_libre(2.0):.2f} m")

# Sur la Lune (g = 1.62 m/s²)
print(f"Lune  : y = {chute_libre(2.0, g=1.62):.2f} m")

# Avec vitesse initiale
print(f"Avec v0=5 : y = {chute_libre(2.0, v0=5):.2f} m")
\end{minted}
\end{codeblock}

\begin{codeoutput}
\begin{verbatim}
Terre : y = 19.62 m
Lune  : y = 3.24 m
Avec v0=5 : y = 29.62 m
\end{verbatim}
\end{codeoutput}

\begin{bonnepratique}
Utilisez des arguments nomm\'es pour am\'eliorer la lisibilit\'e~:
\py{chute\_libre(t=2.0, g=1.62)} est plus clair que \py{chute\_libre(2.0, 1.62)}.
\end{bonnepratique}

% ============================================================
\section{Fonctions avec plusieurs valeurs de retour}

\begin{codeblock}[title=\textbf{Python}]
\begin{minted}{python}
import math

def coordonnees_polaires(x, y):
    """Convertit des coordonnées cartésiennes en polaires."""
    r = math.sqrt(x**2 + y**2)
    theta = math.atan2(y, x)
    return r, theta    # retourne un tuple

r, angle = coordonnees_polaires(3, 4)
print(f"r = {r:.2f}, θ = {math.degrees(angle):.1f}°")
\end{minted}
\end{codeblock}

\begin{codeoutput}
\begin{verbatim}
r = 5.00, θ = 53.1°
\end{verbatim}
\end{codeoutput}

% ============================================================
\section{Port\'ee des variables~: locale vs globale}
\index{port\'ee}\index{variable locale}\index{variable globale}

\begin{definition}[Port\'ee (\emph{scope})]
Une variable d\'efinie \`a l'int\'erieur d'une fonction est \textbf{locale}~: elle n'existe
que pendant l'ex\'ecution de la fonction. Une variable d\'efinie en dehors est
\textbf{globale}.
\end{definition}

\begin{codeblock}[title=\textbf{Python}]
\begin{minted}{python}
x = 10          # variable globale

def ma_fonction():
    x = 5       # variable locale (différente !)
    print(f"Dans la fonction : x = {x}")

ma_fonction()
print(f"En dehors : x = {x}")
\end{minted}
\end{codeblock}

\begin{codeoutput}
\begin{verbatim}
Dans la fonction : x = 5
En dehors : x = 10
\end{verbatim}
\end{codeoutput}

\begin{attention}
\'Evitez d'utiliser \py{global} pour modifier des variables globales depuis une
fonction. Pr\'ef\'erez \emph{toujours} passer les valeurs en param\`etres et utiliser
\py{return}.
\end{attention}

% ============================================================
\section{Fonctions lambda}
\index{lambda}

Les fonctions \py{lambda} sont des fonctions anonymes courtes, utiles pour des
op\'erations simples.

\begin{codeblock}[title=\textbf{Python}]
\begin{minted}{python}
# Fonction lambda pour convertir Celsius en Kelvin
celsius_to_kelvin = lambda T: T + 273.15

temperatures_C = [-40, 0, 20, 37, 100]
temperatures_K = [celsius_to_kelvin(T) for T in temperatures_C]

for c, k in zip(temperatures_C, temperatures_K):
    print(f"{c:>4}°C = {k:>6.2f} K")
\end{minted}
\end{codeblock}

\begin{codeoutput}
\begin{verbatim}
 -40°C = 233.15 K
   0°C = 273.15 K
  20°C = 293.15 K
  37°C = 310.15 K
 100°C = 373.15 K
\end{verbatim}
\end{codeoutput}

% ============================================================
\section{Erreurs courantes}

\begin{errmark}
\begin{enumerate}
    \item \textbf{Oublier \py{return}}~: sans \py{return}, une fonction renvoie
    \py{None} par d\'efaut. Le calcul est effectu\'e mais le r\'esultat est perdu~!
    \begin{minted}{python}
def aire_cercle(r):
    a = 3.14159 * r**2
    # Oubli de return a !

resultat = aire_cercle(5)
print(resultat)  # Affiche None
    \end{minted}

    \item \textbf{Argument mutable par d\'efaut}~: utiliser une liste comme valeur
    par d\'efaut est dangereux, car elle est partag\'ee entre tous les appels~!
    \begin{minted}{python}
# ERREUR : la liste est partagée
def ajouter(element, liste=[]):
    liste.append(element)
    return liste

print(ajouter(1))  # [1]
print(ajouter(2))  # [1, 2] au lieu de [2] !

# CORRECT : utiliser None
def ajouter(element, liste=None):
    if liste is None:
        liste = []
    liste.append(element)
    return liste
    \end{minted}

    \item \textbf{Confondre \py{print} et \py{return}}~: \py{print} affiche \`a l'\'ecran
    mais ne renvoie pas de valeur utilisable dans une expression.
\end{enumerate}
\end{errmark}

% ============================================================
\section{Mini-projet~: Calculateur de formules physiques}

\begin{projet}
Cr\'eez un module contenant les fonctions suivantes~:
\begin{itemize}
    \item \py{energie\_cinetique(m, v)}~: $E_c = \frac{1}{2}mv^2$
    \item \py{force\_gravitationnelle(m1, m2, r)}~: $F = G\frac{m_1 m_2}{r^2}$
    \item \py{periode\_pendule(L, g=9.81)}~: $T = 2\pi\sqrt{L/g}$
    \item \py{loi\_ohm(V=None, R=None, I=None)}~: $V = RI$ (donner deux grandeurs, calculer la troisi\`eme)
\end{itemize}
Chaque fonction doit avoir une docstring compl\`ete. Testez-les avec des valeurs connues.
\end{projet}

\begin{codeblock}[title=\textbf{Python}]
\begin{minted}{python}
import math

def periode_pendule(L, g=9.81):
    """Période d'un pendule simple : T = 2π√(L/g)."""
    return 2 * math.pi * math.sqrt(L / g)

def loi_ohm(V=None, R=None, I=None):
    """Loi d'Ohm : V = R × I. Fournir deux grandeurs."""
    if V is None:
        return R * I
    elif R is None:
        return V / I
    elif I is None:
        return V / R

# Tests
print(f"Pendule L=1m : T = {periode_pendule(1.0):.3f} s")
print(f"Loi d'Ohm : V = {loi_ohm(R=100, I=0.5):.1f} V")
print(f"Loi d'Ohm : I = {loi_ohm(V=12, R=100):.3f} A")
\end{minted}
\end{codeblock}

\begin{codeoutput}
\begin{verbatim}
Pendule L=1m : T = 2.006 s
Loi d'Ohm : V = 50.0 V
Loi d'Ohm : I = 0.120 A
\end{verbatim}
\end{codeoutput}

% ============================================================
\section{Exercices}

\begin{exercice}[$\star$ --- Conversion de temp\'eratures]
\'Ecrivez deux fonctions \py{celsius\_to\_fahrenheit(T)} et \py{fahrenheit\_to\_celsius(T)}.
Testez avec l'eau bouillante (100\,\textdegree C = 212\,\textdegree F).
\end{exercice}

\begin{exercice}[$\star$ --- Factorielle]
\'Ecrivez une fonction \py{factorielle(n)} qui calcule $n!$ avec une boucle.
V\'erifiez que \py{factorielle(10) == 3628800}.
\end{exercice}

\begin{exercice}[$\star\star$ --- Fonction r\'ecursive]
R\'e\'ecrivez la factorielle de mani\`ere \emph{r\'ecursive}\index{r\'ecursion}~: une fonction qui
s'appelle elle-m\^eme. Quel est le cas de base~?
\end{exercice}

\begin{exercice}[$\star\star$ --- Statistiques descriptives]
\'Ecrivez une fonction \py{statistiques(donnees)} qui re\c{c}oit une liste de nombres
et renvoie la moyenne, la m\'ediane et l'\'ecart-type (sans utiliser de biblioth\`eque).
\end{exercice}

\begin{exercice}[$\star\star\star$ --- M\'ethode de Newton]
Impl\'ementez la m\'ethode de Newton\index{m\'ethode de Newton} pour trouver les z\'eros
d'une fonction~: $x_{n+1} = x_n - \frac{f(x_n)}{f'(x_n)}$.
Votre fonction prendra en param\`etres \py{f}, \py{df} (d\'eriv\'ee), \py{x0}
(estimation initiale) et \py{tol} (tol\'erance). Testez avec
$f(x) = x^2 - 2$ pour retrouver $\sqrt{2}$.
\end{exercice}

% ============================================================
\begin{formules}{R\'esum\'e du chapitre~: Fonctions et port\'ee}
\begin{itemize}
    \item \textbf{D\'efinition}~: \py{def nom(params):} ... \py{return resultat}
    \item \textbf{Docstring}~: cha\^ine de documentation entre triple guillemets
    \item \textbf{Arguments par d\'efaut}~: \py{def f(x, n=10):} --- n vaut 10 si non pr\'ecis\'e
    \item \textbf{Arguments nomm\'es}~: \py{f(x=3, n=20)} pour plus de clart\'e
    \item \textbf{Valeurs multiples}~: \py{return a, b} renvoie un tuple
    \item \textbf{Port\'ee locale}~: les variables dans une fonction sont isol\'ees
    \item \textbf{Lambda}~: \py{lambda x: x**2} pour des fonctions courtes
    \item \textbf{Principe DRY}~: ne jamais copier-coller du code, \'ecrire une fonction
\end{itemize}
\end{formules}
